% Part: normal-modal-logic
% Chapter: frame-correspondence
% Section: properties-accessibility

\documentclass[../../../include/open-logic-section]{subfiles}

\begin{document}

\olfileid{nml}{frd}{acc}

\olsection{प्राप्यता संबंधांचे गुणधर्म}

अनेक मोडल !!{formula}s प्राप्यता संबंधाच्या साध्या, कधीकधी
परिचित, गुणधर्मांचे वैशिष्ट्य दर्शवतात. एका दिशेने याचा
अर्थ असा की, एखादा गुणधर्म असलेल्या प्रत्येक प्रतिमानात
संबंधित !!{formula} आणि त्याचे सर्व आदेशन-दृष्टांत सत्य
असतात. प्राप्यता संबंधांचे पाच अभिजात प्रकार आणि त्यांनी
सत्य ठरवलेली सूत्रे प्रथम पाहू.

\begin{thm}\ollabel{thm:soundschemas}
  $\mModel{M} = \tuple{W, R, V}$ हे प्रतिमान असू द्या.
  $R$ ला \olref{tab:five} च्या डाव्या बाजूचा गुणधर्म
  असेल, तर उजव्या बाजूच्या !!{formula} चा प्रत्येक
  दृष्टांत~$\mModel{M}$ मध्ये सत्य असतो.
\end{thm}

\begin{table}[t]
    \begin{tabular}{| l || l |}
      \hline
      {\emph{$R$ पुढील प्रकारचा असेल \dots}} & {\emph{तर \dots सूत्र~$\mModel{M}$ मध्ये सत्य:}} \\
      \hline \hline
      \emph{सर्वत्र उत्तरवर्ती}: $\forall u \exists v Ruv$ & \hbox to.43\textwidth
           {$\Box p \lif \Diamond p$ \hfill (\Ax{D})} \\
      \hline
      \emph{परावर्ती}: $\forall w Rww$
      & \hbox to.43\textwidth
          {$\Box p \lif p$ \hfill (\Ax{T})} \\
      \hline
      \emph{सममित}: &  \hbox to .43\textwidth
           {$p \lif \Box\Diamond p$ \hfill (\Ax{B})} \\
           $\forall u\forall v(Ruv \lif Rvu)$ & \\
      \hline
      \emph{संक्रमक}: & \hbox to.43\textwidth
           {$\Box p \lif \Box \Box p$ \hfill (\Ax{4})} \\
      $\forall u \forall v \forall w ((Ruv \land Rvw) \lif Ruw)$ & \\
      \hline
      \emph{युक्लिडी}: & \hbox to .43\textwidth
        {$\Diamond p \lif \Box\Diamond p$ \hfill (\Ax{5})} \\
      $\forall w \forall u \forall v ((Rwu \land Rwv) \lif Ruv)$ &\\
      \hline
    \end{tabular}
    \caption{अनुरूपतेविषयी पाच तथ्ये.}
    \ollabel{tab:five}
  \end{table}


\begin{proof}
  \Ax{B} चे प्रकरण पाहू. रूपबंध प्रतिमानात सत्य आहे हे
  दाखवण्यासाठी त्याचे सर्व दृष्टांत त्या प्रतिमानातील
  प्रत्येक जगात सत्य आहेत हे दाखवावे लागते. म्हणून
  $!A \lif \Box\Diamond !A$ हा \Ax{B} चा कोणताही
  दृष्टांत, आणि $w \in W$ हे कोणतेही जग असू द्या.
  $\Box \Diamond
  !A$ हे $w$ येथे सत्य असल्याचे दाखवण्यासाठी पूर्वांग
  $!A$ हे $w$ येथे सत्य आहे असे समजा. मग $w$ पासून
  प्राप्य असलेल्या प्रत्येक $w'$ येथे $\Diamond !A$
  सत्य आहे हे दाखवायचे आहे. आता $Rww'$ असलेल्या
  कोणत्याही $w'$ साठी, सममितीच्या गृहीतकावरून $Rw'w$
  देखील असते (\olref{fig:Bsymm} पाहा).
  $\mSat{M}{!A}[w]$ असल्यामुळे
  $\mSat{M}{\Diamond !A}[w']$. $Rww'$ असलेले $w'$
  कोणतेही असल्यामुळे $\mSat{M}{\Box\Diamond!A}[w]$.

  इतर प्रकरणे सरावासाठी ठेवली आहेत.
\end{proof}

\begin{prob}
  \olref[nml][frd][acc]{thm:soundschemas} ची सिद्धता पूर्ण करा.
\end{prob}

\begin{figure}
  \begin{center}
    \begin{tikzpicture}[modal]
      \node[world] (w1) [label={below:$\mSat{{}}{\formula{A}}$\\
          $\mSat{{}}{\Box\Diamond \formula{A}}$}] {$w$};
      \node[world] (w2) [label=below:{$\mSat{{}}{\Diamond \formula{A}}$},
        right=of w1] {$w'$};
      \draw[->, bend left] (w1) to (w2);
      \draw[->, bend left] (w2) to (w1);
    \end{tikzpicture}
  \end{center}
\caption{सममितीवर आधारित युक्तिवाद.}
\ollabel{fig:Bsymm}
\end{figure}

\olref{thm:soundschemas} मधील व्यत्यास खरे नाहीत, हे लक्षात
घ्या: एखादा रूपबंध प्रतिमानात सत्य आहे म्हणून त्या प्रतिमानाच्या
प्राप्यता संबंधाला संबंधित गुणधर्म असतोच असे नाही.
\Ax{T} आणि परावर्ती प्रतिमाने या प्रकरणात, \Ax{T} स्वतःच
असत्य असलेले प्रतिमान सहज देता येते: $W = \{w\}$ आणि
$V(p) = \emptyset$ घ्या; प्राप्यता संबंध रिकामा असू द्या.
मग $R$ परावर्ती नाही, पण $\mSat{M}{\Box
  p}[w]$ आणि $\mSat/{M}{p}[w]$. तरी येथे~$\mModel{M}$
मध्ये \Ax{T} चा एकच दृष्टांत असत्य आहे; इतर दृष्टांत,
उदाहरणार्थ $\Box \lnot p \lif \lnot p$, सत्य आहेत.
\Ax{T} चे \emph{प्रत्येक आदेशन-दृष्टांत}~$\mModel{M}$
मध्ये सत्य असूनही $\mModel{M}$ परावर्ती नसते, असे
प्रतिमान शोधणे अधिक कठीण आहे. पण अशी प्रतिमानेही आहेत:

\begin{prop}\ollabel{prop:reflexive}
  $\mModel{M} = \tuple{W, R, V}$ हे प्रतिमान असू द्या,
  जिथे $W = \{u, v
  \}$, जगं $u$ आणि $v$ यांचा $R$ संबंध दोन्ही दिशांनी
  आहे, म्हणजे $Ruv$ आणि $Rvu$, आणि कोणतेही जग स्वतःशी
  संबंधित नाही. प्रत्येक $p$ साठी $u \in V(p) \Leftrightarrow v
  \in V(p)$ असे समजा. मग:
  \begin{enumerate}
  \item प्रत्येक $!A$ साठी: $\mSat{M}{!A}[u]$ तेव्हा आणि
    केवळ तेव्हाच, जेव्हा $\mSat{M}{!A}[v]$
    ($!A$ वर विगमन वापरा).
  \item \Ax{T} चा प्रत्येक दृष्टांत $\mModel{M}$ मध्ये सत्य आहे.
  \end{enumerate}
  $\mModel{M}$ परावर्ती नाही (खरेतर ते
  \emph{अपरावर्ती} आहे), म्हणून \Ax{T} च्या बाबतीत
  \olref{thm:soundschemas} चा व्यत्यास असत्य आहे.
  \olref{thm:soundschemas} मधील काही, पण सर्वच नव्हे,
  इतर रूपबंधांसाठीही असेच युक्तिवाद देता येतात.
\end{prop}

\begin{prob}
  \olref[nml][frd][acc]{prop:reflexive} मधील दावे सिद्ध करा.
\end{prob}

आपण \Ax{D}, \Ax{T}, \Ax{B}, \Ax{4}, आणि \Ax{5}
या पाच अभिजात !!{formula}s वर भर देणार असलो, तरी
\olref{tab:anotherfive} मध्ये प्राप्यता संबंधांचे आणखी
काही गुणधर्म नोंदवले आहेत. प्रत्येक जगापासून जास्तीत
जास्त एक जग प्राप्य असेल, तर प्राप्यता संबंध~$R$
आंशिक फलनीय असतो. प्रत्येक जगापासून नेमके
एकच जग प्राप्य असेल, तर त्याला फलनीय म्हणतो.
(म्हणून फलनीय संबंध नेमके तेच आहेत जे सर्वत्र उत्तरवर्ती
आणि आंशिक फलनीय दोन्ही आहेत.) त्यांना ``फलनीय''
म्हणतात कारण प्राप्यता संबंध एखाद्या (आंशिक) फलनासारखा
वागतो. $Ruv$ असेल तेव्हा $u$ आणि~$v$ यांच्या
``मध्ये'' एखादे~$w$ असेल, तर संबंध दुर्बल सघन
असतो. एका अर्थाने दुर्बल सघन संबंध संक्रमक संबंधांच्या
उलट आहेत: संक्रमक संबंधात~$u$ पासून~$w$ मार्गे
$v$ पर्यंत जाता येत असेल, तर $u$ पासून $v$ पर्यंत
थेट जाता येते; दुर्बल सघन संबंधात~$u$ पासून $v$
पर्यंत थेट जाता येत असेल, तर एखाद्या~$w$ मार्गेही
जाता येते. एखाद्या जग~$w$ पासून जगं $u$ आणि~$v$
प्राप्य असतील तेव्हा $u$ आणि~$v$ या दोन्हींपासून
एकच जग~$t$ प्राप्य असेल, तर संबंध दुर्बल दिशित असतो;
यालाच कधीकधी ``हिऱ्याचा गुणधर्म'' किंवा ``संगमन''
म्हणतात.

\begin{table}[t]
    \begin{tabular}{| p{.50\textwidth} || p{.43\textwidth} |}
      \hline
      {\emph{$R$ पुढील प्रकारचा असेल \dots}} & {\emph{तर \dots सूत्र~$\mModel{M}$ मध्ये सत्य:}} \\
      \hline \hline
      \emph{आंशिक फलनीय}: &
      \multirow{2}{*}{$\Diamond p \lif \Box p$} \\
      $ \forall w \forall u \forall v ((Rwu \land Rwv) \lif u =v )$ & \\
      \hline
      \multirow{2}{*}{\emph{फलनीय}: $\forall w \exists v \forall u(Rwu \liff \eq[u][v]) $}
      & \multirow{2}{*}{$\Diamond p \liff \Box p$} \\
      & \\
      \hline
      \emph{दुर्बल सघन}: &  \multirow{2}{*}{$\Box \Box p \lif
        \Box p$} \\
      $\forall u\forall v(Ruv \lif \exists w(Ruw \land Rwv))$ & \\
      \hline
      \emph{दुर्बल संयोजित}: &
      \multirow{3}{*}{\hbox to.43\textwidth{$
        \begin{array}{@{}l@{}}
          \Box ((p \land \Box p) \lif q) \lor {} \\
          \qquad \Box ((q \land \Box q) \lif p)
        \end{array}$ \hfill (\Ax{L})}
      } \\
      $\forall w \forall u \forall v ((Rwu \land Rwv) \lif {}$ &\\
      \qquad $(Ruv \lor u=v \lor Rvu))$ & \\
      \hline
      \emph{दुर्बल दिशित}: & \multirow{3}{*}{\hbox to .43\textwidth
        {$\Diamond\Box p \lif \Box\Diamond p$\hfill (\Ax{G})}} \\
      $\forall w \forall u \forall v ((Rwu \land Rwv) \lif {}$ & \\
      \qquad $\exists t (Rut \land Rvt))$ & \\
      \hline
    \end{tabular}
    \caption{अनुरूपतेविषयी आणखी पाच तथ्ये.}
    \ollabel{tab:anotherfive}
  \end{table}

\begin{prob}
  $\mModel{M} = \tuple{W, R, V}$ हे प्रतिमान असू द्या.
  $R$ ला \olref[nml][frd][acc]{tab:anotherfive} मधील
  डावीकडील गुणधर्म असेल, तर उजवीकडील संबंधित सूत्राचा
  प्रत्येक दृष्टांत~$\mModel{M}$ मध्ये सत्य असतो, हे दाखवा.
\end{prob}



\end{document}
